\centerline{\bf \Huge Алгебра и начало анализа}
\centerline{\bf \Huge Теорема о промежуточном значении.}

\begin{flushright}
        \scriptsize{}
\end{flushright}

При решении задач о квадратном трёхчлене часто бывает полезно следующее утверждение, которое называют \textbf{теоремой о промежуточном значении для квадратного трёхчлена.} Пусть $f(x)=a x^2+b x+c$, где $a \neq 0$. Если $f(\alpha) f(\beta) \leqslant 0$, то на отрезке $[\alpha, \beta]$ лежит по крайней мере один корень квадратного уравнения $a x^2+b x+c=0$.

Для $f(x)=\varphi(x)-\psi(x)$ это означает, что если, например, $\varphi(\alpha) \geqslant \psi(\alpha)$ и $\varphi(\beta) \leqslant \psi(\beta)$, то на отрезке $[\alpha, \beta]$ есть точка $x_0$, для которой $\varphi\left(x_0\right)=\psi\left(x_0\right)$.

Вообще говоря, это всё лишь частный случай одной большой теоремы Больцана-Коши из математического анализа, которая говорит, что если непрерывная функция, определённая на вещественном промежутке, принимает два значения, то она принимает и любое значение между ними.

\problem{0}
Докажите теорему о промежуточном значении для квадратного трёхчлена.

\problem{1}
Пусть $\alpha$ - корень уравнения $x^2+a x+b=0, \beta$ - корень уравнения $x^2-a x-b=0$. Докажите, что между числами $\alpha$ и $\beta$ есть корень уравнения $x^2-2 a x-2 b=0$.


\problem{2}
Квадратный трёхчлен $f(x)$ имеет два вещественных корня, разность которых не меньше натурального числа $n \geqslant$ $\geqslant 2$. Докажите, что квадратный трёхчлен $f(x)+f(x+1)+\ldots$ $\ldots+f(x+n)$ имеет два вещественных корня.

\problem{3}
Даны уравнения $a x^2+b x+c=0$ и $-a x^2+b x+c=0$. Докажите, что если $x_1$ - корень первого уравнения, а $x_2$ второго, то найдётся корень $x_3$ уравнения $\frac{a}{2} x^2+b x+c=0$, для которого либо $x_1 \leqslant x_3 \leqslant x_2$, либо $x_1 \geqslant x_3 \geqslant x_2$.

\problem{4}
Докажите, что на отрезке $[-1,1]$ квадратный трёхчлен $f(x)=x^2+a x+b$ принимает значение, абсолютная величина которого не меньше $1 / 2$.

\problem{5}
Докажите, что если $\left|a x^2+b x+c\right| \leqslant 1 / 2$ при всех $|x| \leqslant 1$, то $\left|a x^2+b x+c\right| \leqslant x^2-1 / 2$ при всех $|x| \geqslant 1$.

\problem{6}
Докажите, что если $\left|a x^2+b x+c\right| \leqslant 1$ при $|x| \leqslant 1$, то $\left|c x^2+b x+a\right| \leqslant 2$ при $|x| \leqslant 1$.